\bookchapter{Penutup: Tetap Ingin Tahu, Tetap Jernih}

Selamat Anda sudah membaca sampai sini.

Buku ini dari hardware ponsel bedah sampai ekosistem software, dari pengaturan harian gali sampai developer mode, dari security protection obrolan sampai warisan digital. Jika membantu Anda selesaikan satu masalah aktual — bahkan jika hanya belajar mematikan satu permission yang tidak seharusnya aktif, membentuk satu kebiasaan backup — maka kata-kata ini tidak sia-sia.

Tapi lebih penting dari skill konkret apa pun, adalah dua hal:

\begin{itemize}
  \item \textbf{Tetap curious}: Ponsel berubah setiap tahun, AI berubah, folding screen berubah, satelit communication berubah. Langkah operasi yang dipelajari hari ini besok mungkin sudah outdated, tapi rasa ingin tahu "ingin tahu kenapa" tidak akan pernah outdated. Tanya satu kalimat lagi "permission ini untuk apa" "data ini disimpan di mana" "fitur ini apakah akan bocor privasi", lebih berguna dari mengingat 100 langkah operasi.
  \item \textbf{Tetap sadar}: Ponsel adalah tool, bukan tuan. Dirancang untuk menarik perhatian Anda, menempati waktu Anda, memandu konsumsi Anda. Memakainya adalah pilihan Anda, dipakai olehnya adalah kehilangan kendali. Secara berkala angkat kepala lihat screen time, secara berkala pikirkan "saya sedang pakai ponsel untuk melakukan sesuatu, atau sedang dihabiskan oleh ponsel", ini bukan dramatis, ini adalah instinct bertahan hidup era digital.
\end{itemize}

Teknologi berubah, tapi kebutuhan manusia tidak berubah — kita butuh rasa aman, butuh koneksi, butuh dihormati. Ponsel hanyalah carrier. “Kelas ponsel” yang sebenarnya tidak ada di buku, ada di momen Anda unlock screen setiap hari.

\begin{tcbpractice}{Latihan Terakhir}
Tutup buku ini, ambil ponsel Anda, lakukan satu hal:\textbf{Pakai 30 detik pikirkan, di ponsel Anda ada tiga hal apa yang “hilang akan benar-benar sakit hati”?}Lalu, pastikan mereka sudah di-backup. Buku ini sudah berharga.

\end{tcbpractice}
